% Generated draft: requires mathematical and source review.
\documentclass[11pt]{article}
\usepackage[margin=25mm]{geometry}
\usepackage{fontspec}
\setmainfont{Latin Modern Roman}
\usepackage{amsmath,amssymb,mathtools,unicode-math}
\setmathfont{Latin Modern Math}
\usepackage{enumitem,needspace,xurl,hyperref}
\setlength{\parindent}{0pt}
\setlength{\parskip}{5pt}
\setlength{\emergencystretch}{3em}
\newcommand{\N}{\mathbb{N}}
\newcommand{\Z}{\mathbb{Z}}
\newcommand{\Q}{\mathbb{Q}}
\newcommand{\R}{\mathbb{R}}
\newcommand{\C}{\mathbb{C}}
\newcommand{\F}{\mathbb{F}}
\newcommand{\A}{\mathbb{A}}
\newcommand{\PP}{\mathbb{P}}
\newcommand{\E}{\mathbb{E}}
\newcommand{\eps}{\varepsilon}
\newcommand{\dd}{\,\mathrm{d}}
\begin{document}
\section*{UnsolvedMath ID: 3100005}
\textbf{Source problem code:} catalogue entry
\subsection{Definitions and mathematical statement}
For a finite set $S\subset\mathbb R^2$, let $\mathcal L(S)$ be the set of distinct lines determined by pairs of distinct points of $S$. In the Euclidean interpretation define
\[
T(S)=\{x\in\mathbb R^2:\text{there exist distinct nonparallel }
\ell,\ell'\in\mathcal L(S)\text{ with }x=\ell\cap\ell'\}.
\]
Coincident lines are one member of $\mathcal L(S)$, and parallel lines add no point at infinity. Starting with $S_0=S$, set $S_{i+1}=T(S_i)$ and $N_i=|S_i|$ for integers $i\geq0$.

\textbf{Problem.} Determine how rapidly $N_i$ grows under repeated line-and-intersection construction, including the dependence on the initial configuration $S$. The source does not specify the size or nondegeneracy of $S$, whether it seeks universal bounds, a typical-configuration rate, or a classification of possible rates. Finite Euclidean sets are the explicit conventional setting above; an infinite or projective-plane version would be a different question.

The operator is the source's replacement $T(S)$, not automatically the accumulated closure $S\cup T(S)$. For degenerate initial sets those definitions can differ, so initial-point retention must not be added silently. No general-position assumption is imposed.

The supplied research summary describes doubly exponential growth as a candidate maximal scale and says the general exact growth question remains unresolved. It mentions rank/lattice-degree arguments, Sylvester--Gallai configurations, and degeneracies only as context, without a precise bound or theorem to incorporate.
\subsection{Short English statement}
Study the possible growth rates obtained by repeatedly intersecting all lines determined by the current finite planar point set. The source leaves the initial-position assumptions and the requested type of growth bound unspecified.
\subsection{Sources}
\begin{itemize}
\item Cooper, \emph{Combinatorial Problems I Like}, Geometry, source-order bullet 5, October 2020 snapshot, \url{https://people.math.sc.edu/cooper/combprob.html}; HTML extraction, accessed 2026-07-29. The author page labels the item open in that snapshot; later resolutions were not checked item by item according to the record.
\item The statement has a see-this link whose target is absent from the extraction. The research summary names Magnuson and collaborators and uses the phrase Hetzheim-type growth, but supplies no identifiable publication, theorem, or discussion URL.
\item Local registry \texttt{data/unsolvedmath/problems.json}, permanent ID $3100005$, source code catalogue entry, \emph{Suppose we begin with a set of points S in the plane}. The record labels this open and OPEN-TRIAGE, with NEEDS\_REVIEW source-status and rights qualifications.
\end{itemize}
\end{document}
